\subsection{SECTION OF AN OPEN SET; SECTION OF A CLOSED SET; PROJECTION OF AN OPEN SET. PARTIAL CONTINUITY}

PROPOSITION 4. Let \(\mathbf{X}_1, \mathbf{X}_2\) be two topological spaces; then for each \(a_1 \in \mathbf{X}_1\) , the mapping \(x_2 \to (a_1, x_2)\) is a homeomorphism of \(\mathbf{X}_2\) onto the subspace \(\{a_1\} \times \mathbf{X}_2\) of \(\mathbf{X}_1 \times \mathbf{X}_2\) .

This is a particular case of Corollary 1 of Proposition 1 of no. 1 applied to the constant function \(x_{2} \rightarrow a_{1}\) .

The mapping \(x_{2} \to (a_{1}, x_{2})\) is a continuous section (§ 3, no. 5) with respect to the equivalence relation \(\mathrm{pr}_{2}z = \mathrm{pr}_{2}z'\) in \(\mathbf{X}_{1} \times \mathbf{X}_{2}\) ; the quotient space \(\mathbf{X}_{1} \times \mathbf{X}_{2}\) by this equivalence relation is therefore homeomorphic to \(\mathbf{X}_{2}\) .

COROLLARY. The section \(\mathbf{A}(x_1)\) of an open (resp. closed) set \(\mathbf{A}\) of the product \(\mathbf{X}_1 \times \mathbf{X}_2\) at an arbitrary point \(x_1 \in \mathbf{X}_1\) (*) is open (resp. closed) in \(\mathbf{X}_2\) .

PROPOSITION 5. The projection of an open set U of the product \(X_{1} \times X_{2}\) onto either factor is an open set.

For example, we have \(\mathrm{pr}_2\mathrm{U} = \bigcup_{x_1\in X_1}\mathrm{U}(x_1)\) , and the proposition follows from the Corollary to Proposition 4 and axiom \((\mathbf{O}_{\mathbf{I}})\) .

\pandocbounded{\includegraphics[keepaspectratio]{images/part001_fc188d82.jpg}}

Remark 1. The projection of a closed subset of \(\mathbf{X}_1 \times \mathbf{X}_2\) onto a factor need not be closed. For example, in the rational plane \(\mathbf{Q}^2\) , the hyperbola whose equation is \(x_1 x_2 = 1\) is a closed set, but both its projections are equal to the complement of the point 0 in \(\mathbf{Q}\) , and this is not a closed set.

PROPOSITION 6. Let \(\mathbf{X}_1, \mathbf{X}_2, \mathbf{Y}\) be three topological spaces, \(f\) a mapping of the product space \(\mathbf{X}_1 \times \mathbf{X}_2\) into \(\mathbf{Y}\) . If \(f\) is continuous at the point \((a_1, a_2)\) then the partial mapping \(x_2 \to f(a_1, x_2)\) of \(\mathbf{X}_2\) into \(\mathbf{Y}\) is continuous at the point \(a_2\) .

For this mapping is the composition of \(f\) and the mapping \(x_{2} \to (a_{1}, x_{2})\) ; hence the result follows from Proposition 4.

Proposition 6 is often expressed by saying that a continuous function of two variables is continuous with respect to each of them separately.

\pandocbounded{\includegraphics[keepaspectratio]{images/part001_31ca3096.jpg}}

Remark 2. It is possible for all the partial mappings determined by a map \(f: \mathbf{X}_{1} \times \mathbf{X}_{2} \to \mathbf{Y}\) to be continuous without \(f\) being continuous on \(\mathbf{X}_{1} \times \mathbf{X}_{2}\) (cf.~Chapter IX, § 5, Exercise 23). * For example if \(f\) is the mapping of the real plane \(\mathbb{R}^{2}\) into \(\mathbb{R}\) defined by

\[
f (x, y) = x y / \left(x ^ {2} + y ^ {2}\right) \quad \text { if } \quad (x, y) \neq (0, 0)
\]

and \(f(0, 0) = 0\) , then all the partial mappings are continuous; but \(f\) is not continuous at \((0, 0)\) , since \(f(x, x) = 1/2\) if \(x \neq 0\) .

If \(g\) is a mapping of \(\mathbf{X}_1\) into \(\mathbf{Y}\) , continuous at a point \(a_1\) , then the mapping \((x_1, x_2) \to g(x_1)\) of \(\mathbf{X}_1 \times \mathbf{X}_1 \to \mathbf{Y}\) is continuous at all points \((a_1, x_2)\) , for it is the composition of \(g\) and the projection onto \(\mathbf{X}_1\) .

The results of this subsection are easily extended to an arbitrary product \(\prod_{i\in I}X_i\) of topological spaces by remarking that this product is homeomorphic to the product \(\left(\prod_{i\in J}X_i\right)\times\left(\prod_{i\in K}X_i\right)\) for any partition (J, K) of I (no. 1, Proposition 2).

